% ============================================================
% conclusion.tex
% ============================================================

\section{Conclusion}
\label{sec:conclusion}

This study examined the relationships between cortical morphometry and four variables of interest, namely sleep quality, fluid intelligence, neuroticism, and age, using two large publicly available datasets that together span the adult lifespan from 22 to 89 years.
Three findings stand out.
First, fluid intelligence was positively associated with cortical surface area across every region tested, confirming that the structural substrate of cognitive ability is broadly distributed rather than confined to a few specialized areas.
Second, aging was the single most powerful determinant of cortical structure, with primary motor and somatosensory cortices showing correlations with age as large as \textit{r}\,=\,$-$.69, underscoring the need for age-appropriate study designs in structural neuroimaging.
Third, and most importantly, neuroticism showed a developmental dissociation: it was associated with reduced surface area (but not thickness) in young adults, and with widespread cortical thinning in older adults, suggesting that the structural consequences of this personality trait accumulate over decades through stress-related biological pathways.

Taken together, these findings demonstrate the value of multi-cohort, multi-metric designs for understanding brain-behavior relationships and highlight the aging cortex as a substrate that is differentially shaped by personality-linked physiological processes over the lifespan.
